
\documentclass[12pt]{article}
\pagestyle{empty} \topmargin -0.5in \textheight 9in \textwidth 6.5in
\oddsidemargin -0in
\usepackage{amssymb}
\usepackage{amsmath}
\usepackage[normalem]{ulem}
\usepackage{enumerate}
\usepackage{xcolor}

\newcommand{\mcE}{\mathcal{E}}

\DeclareMathOperator{\tr}{tr}

\newcommand{\vect}[1]{\mathbf{#1}}

\begin{document}

\begin{flushleft}

\textcolor{blue}{\textbf{NAME: Andrew Yonan}}\\
\textbf{Department of Applied Mathematics and Statistics}

\textbf{COLORADO SCHOOL OF MINES}\\


MATH $455$: Partial Differential Equations
\end{flushleft}

\vspace{0.1in}

\begin{center}
{\bf Assignment \#6 - Problems in Higher Dimensions} \\
Due Tuesday, November 17, 2026 \\
\end{center}

\vspace{0.2in}


\noindent $\textbf{1}$. (12 points) Consider the doubly-indexed Fourier series
$$f(x,y) = \sum_{m=1}^\infty \sum_{n=1}^\infty b_{mn} \sin \left ( \frac{n\pi x}{a} \right )\sin \left ( \frac{m\pi y}{b} \right )$$
for a given function $f \in \mathrm{PC}^1 \left ( [0,a] \times [0,b] \right ).$
\begin{enumerate}[(a)]
\item Rewrite the series in the form
\begin{equation}
\label{f}
f(x,y) = \sum_{m=1}^\infty B_m(x) \sin \left ( \frac{m\pi y}{b} \right )
\end{equation}
and identify $B_m(x)$ in terms of the expressions given above. \\ \\
\textcolor{blue}{\textbf{Solution} — 
\begin{align*}
f(x,y) =& \sum_{m=1}^\infty \sum_{n=1}^\infty b_{mn} \sin \left ( \frac{n\pi x}{a} \right )\sin \left ( \frac{m\pi y}{b} \right ) \\
=& \sum_{m=1}^\infty \sin \left ( \frac{m\pi y}{b} \right ) \underbrace{ \left[  \sum_{n=1}^\infty b_{mn} \sin \left ( \frac{n\pi x}{a} \right ) \right] }_{B_m(x)}
\end{align*}
$$\boxed{B_m(x) =  \sum_{n=1}^\infty b_{mn} \sin \left ( \frac{n\pi x}{a} \right )}$$
}

\item Consider $0 < x < a$ fixed. Then \eqref{f} is a Fourier series in $y$. Using our previous knowledge of (singly-indexed) Fourier series, write down an expression for $B_m(x)$ in terms of $f(x,y)$ and $\sin \left ( \frac{m\pi y}{b} \right )$. \\ \\
\textcolor{blue}{\textbf{Solution } — If we fix $x = x_0$, then $B_m(x) = B_m(x_0)$ becomes a function only of $m$. Let us call this function $C_m = B_m(x_0)$. So
\begin{align*}
f(x_0,y) &= \sum_{m=1}^\infty \sin \left ( \frac{m\pi y}{b} \right ) B_m(x_0) \\
&= \sum_{m=1}^\infty C_m \sin \left ( \frac{m\pi y}{b} \right )
\end{align*}
Since this is the Fourier sin-expansion of $f(x_0, y)$ on $y \in [0, b]$, then the Fourier coefficient $C_m$ must be
\begin{align*}
C_m = \frac{2}{b} \int_0^b f(x_0, y) \sin \left( \frac{m \pi y}{b} \right) dy
\end{align*}
So
\begin{align*}
\boxed{B_m(x) = \frac{2}{b} \int_0^b f(x, y) \sin \left( \frac{m \pi y}{b} \right) dy}
\end{align*}
}

\item Consider $m \in \mathbb{N}$ fixed. Using the expression for $B_m(x)$ derived in part (a) and our previous knowledge of Fourier series, solve for $b_{mn}$ in terms of $B_m(x)$ and $\sin \left ( \frac{n\pi x}{a} \right )$. \\ \\
\textcolor{blue}{\textbf{Solution} — Fix $m = m_0$. Then from
\begin{align*}
B_m(x) &=  \sum_{n=1}^\infty b_{mn} \sin \left ( \frac{n\pi x}{a} \right ) \\
\implies& B_{m_0}(x) =  \sum_{n=1}^\infty b_{m_0n}\sin \left ( \frac{n\pi x}{a} \right )
\end{align*}
Hence 
\begin{align*}
B_{m_0n} = \frac{2}{a} \int_0^a B_{m_0}(x) \sin \left( \frac{n \pi x}{a} \right) dx \\
\end{align*}
And, re-generalizing to to non-fixed $m$, 
\begin{align*}
\boxed{b_{mn} = \frac{2}{a} \int_0^a B_{m}(x) \sin \left( \frac{n \pi x}{a} \right) dx} \\
\end{align*}
}

\item Finally, substitute the formula for $B_m$ derived in part (b) within the formula for $b_{mn}$ found by part (c), and use this to identify a final formula for the Fourier coefficients $b_{mn}$ that depends upon $f(x,y)$. \\  \\
\textcolor{blue}{\textbf{Solution} — 
\begin{align*}
b_{mn} &= \frac{2}{a} \int_0^a B_{m}(x) \sin \left( \frac{n \pi x}{a} \right) dx \\
&= \frac{2}{a} \int_0^a \left[ \frac{2}{b} \int_0^b f(x, y) \sin \left( \frac{m \pi y}{b} \right) dy \right] \sin \left( \frac{n \pi x}{a} \right) dx \\
&= \boxed{\frac{4}{ab} \int_0^a  \int_0^b f(x, y) \sin \left( \frac{n \pi x}{a} \right) \sin \left( \frac{m \pi y}{b} \right) dy dx} \\
\end{align*}
}

\end{enumerate}


\vspace{0.3in}

\noindent $\textbf{2}$. (18 points) 
Find the solution $u(t,x,y)$ of the wave equation
\begin{equation}
\label{WE}
\tag{WE}
\frac{\partial^2 u}{\partial x^2} + \frac{\partial^2 u}{\partial y^2} - \frac{1}{c^2} \frac{\partial^2 u}{\partial t^2}  = 0
\end{equation}
on the square 
$$\Omega = \{(x,y) : 0 < x < \pi , 0 < y < \pi \}$$
with homogenous Dirichlet boundary conditions and the initial conditions 
$$u(0,x,y) = xy(\pi-x)(\pi-y) \qquad \mathrm{and} \qquad \frac{\partial u}{\partial t}(0,x,y) = 0$$
for $(x,y) \in \Omega$.
As in class, you may assume that the solution is oscillatory in $t$. Also, notice the symmetry in $u(0,x, y)$.
Upon separating variables, you do NOT need to go through all of the cases to solve for the eigenfunctions and eigenvalues. Just use our previous work!

\end{document}
